\documentclass[11pt]{article}
\usepackage[margin=1in]{geometry}
\usepackage{booktabs}
\usepackage{hyperref}
\usepackage{url}
\usepackage{parskip}

\title{Replication Audit: Dose-Rate Effect on Fast-Neutron-Induced DSBs in Human Lymphocytes}
\author{LUCID-100 Replication Program (Ollie / Rick Stevens AI)}
\date{2026-07-06 (backfill of 2026-06-22 audit)}

\begin{document}
\maketitle

\section*{Paper}
Nair S., Engelbrecht M., Miles X., Ndimba R., Fisher R., du Plessis P., Bolcaen J.,
Nieto-Camero J., de Kock E., Vandevoorde C. (2019).
\emph{The Impact of Dose Rate on DNA Double-Strand Break Formation and Repair in
Human Lymphocytes Exposed to Fast Neutron Irradiation.}
\textbf{Int. J. Mol. Sci.} 20(21):5350.
DOI: \href{https://doi.org/10.3390/ijms20215350}{10.3390/ijms20215350}
(PMID 31661782, PMCID PMC6862539, OA CC-BY 4.0).

\section*{Verdict}
\textbf{REPLICATED} at table-level scope. Coverage 9/10, Agreement 8/10.
Of 8 testable headline claims, 7 verify cleanly; 1 (repair half-life) verifies
numerically but only under a fit variant that contradicts the paper's own stated
Methods text.

\section{Scope}
The paper is a wet-lab radiobiology assay measuring \(\gamma\)-H2AX foci in
peripheral blood lymphocytes (n=4 donors) exposed to iThemba LABS p(66)/Be(40)
fast neutrons at two dose rates (HDR = 0.400~Gy/min, LDR = 0.015~Gy/min;
ratio 26.7$\times$). It reports (a) a dose-response induction curve fit to a
2nd-order polynomial, (b) per-dose HDR/LDR ratios, (c) a K-coefficient
(foci/Gy) analysis, and (d) a repair-kinetics time-course at 1~Gy fit to a
single-exponential decay.

Replication was table-level only. No raw per-cell foci counts, no
Metafer/MetaCyte classifier settings, no GraphPad Prism v5 project, no
neutron-spectrum file, and no raw micrographs are deposited (EPMC
\texttt{hasSuppl:N}, no Zenodo/Figshare/GEO/GitHub). All wet-lab and image-analysis
layers are therefore permanently blocked from independent re-execution.

\section{Method}
We digitized Tables 1--3 from the OA PDF via \texttt{pdftotext}, implemented
the paper's fits in NumPy (Python 3.14), and ran two harnesses:
\texttt{scripts/smoke\_replicate.py} (3 checks, first-pass) and
\texttt{scripts/extended\_replicate.py} (7 checks, this audit, E1--E7). Total
wall-time $<$ 2~s. No paid endpoints, no author contact, no heavy compute.
The repair half-life fit was performed in three variants: (A) raw
single-exponential from $t\!\ge\!2$~h, (B) the paper's stated method
(subtract $foci(24\text{h})$, drop 24~h, fit $t\!\in\!\{2,4,8,12\}$), and (C)
subtract-and-keep-24h.

\section{Results}
\begin{tabular}{clll}
\toprule
\# & Claim & Paper & Replicated \\
\midrule
C1 & Mean HDR/LDR ratio (aggregate) & +40\,\% & +39.77\,\% \\
C2 & Per-dose HDR/LDR ratios (Table 2) & [1.22,1.87,1.44,1.30,1.16] & max err 0.005 \\
C3 & Max K-coefficient locus & 1.87 @ 0.250~Gy & 1.872 @ 0.250~Gy \\
C4 & Poly2 induction fit quality & (not reported) & $R^2$=0.994/0.988 \\
C5 & Dose-rate ratio & 26.67$\times$ & 26.67$\times$ \\
C6 & 24~h residual difference & n.s. & Welch $p$=0.306 \\
C7 & Repair half-life & 8.6~h HDR / 12.0~h LDR & 9.92/13.08 (var. A); 2.89/3.27 (var. B) \\
C8 & Ordering LDR $>$ HDR half-life & LDR $>$ HDR & LDR $>$ HDR (both variants) \\
\bottomrule
\end{tabular}

\textbf{Score:} 7 fully verified, 1 verified-with-caveat. The C7 anomaly is
structural: the fit method \emph{as literally described in Methods} yields
$\sim$3~h half-lives, not 8.6/12.0~h. Only the raw single-exponential fit
(no subtraction) reproduces the reported numbers within 10--15\%.

\section{Critique (Rick's hard requirement --- genuine, not whitewash)}
\begin{itemize}
  \item \textbf{Fast-neutron RBE not reproduced from data --- assumed.} The paper
        does not report an RBE calculation against a low-LET reference. There is
        no matched $^{60}$Co / X-ray dose-response curve in the same experimental
        campaign, so any RBE inference is implicit / cross-study. Our
        replication cannot reproduce an RBE value because none is stated.
  \item \textbf{DRE at high LET \emph{qualitatively} present but not quantified.} The
        paper's central claim --- a $\sim$40\,\% HDR-over-LDR induction effect ---
        \emph{is} a dose-rate effect (DRE) in the DSB-yield direction, but the
        classical high-LET expectation is that DREs \emph{shrink} relative to low
        LET. The paper does not compute a DRE-reduction factor or compare
        against a low-LET DRE baseline; our replication confirms the effect
        magnitude but cannot compare it to a low-LET control in-paper.
  \item \textbf{Donor variability under-characterised.} $n=4$ donors, and only
        cross-donor mean $\pm$ SD is reported. Per-donor breakdowns are not in
        the paper. With no per-cell CSV we cannot bootstrap per-donor
        variability, so all confidence intervals we produce are optimistic
        (they treat within-donor and between-donor variance as a single pooled
        SD).
  \item \textbf{Dose-response linearity not tested against LQ.} The paper
        chooses poly2 without comparing against a linear or linear-quadratic
        (LQ) $\alpha D + \beta D^2$ model. Our AICc comparison (E4) shows poly2
        beats linear by $\Delta$AICc $\sim$5--6 units, but LQ (no intercept)
        is not distinguishable from poly2 given only 5 dose points --- this
        is a real underpowered-model-selection issue.
  \item \textbf{Neutron spectrum only labelled ``fast'', not fully specified.}
        The paper cites prior dosimetry work (ref [42]) but does not include
        the p(66)/Be(40) energy spectrum, mean/median neutron energy, or LET
        distribution at the irradiation point. Without the spectrum file,
        no mechanistic Monte Carlo (PARTRAC / MCDS / TRAX) DSB-induction
        replication is possible, and cross-facility comparisons collapse.
  \item \textbf{Method-text vs numbers inconsistency (C7).} The stated
        subtract-baseline-and-fit procedure produces 2.89/3.27~h half-lives,
        not 8.6/12.0~h. The reported numbers are recoverable only under a raw
        single-exponential fit or (equivalently) a fit that treats the
        baseline as a fitted plateau rather than a subtracted constant. This
        is a real reproducibility hazard independent of data availability.
\end{itemize}

\section{Repro Ceiling}
Table-level reproducibility is high (100\,\% of testable headline claims tested,
87.5\,\% cleanly pass). Deeper reproducibility --- per-donor stats, ANOVA,
image analysis, mechanistic modelling --- is blocked by six missing artifacts:
per-cell $\gamma$-H2AX foci counts, Metafer/MetaCyte classifier config,
GraphPad Prism v5 project, neutron energy spectrum, raw micrographs, and
IRB/consent chain. All six are documented in \S7 of the underlying REPORT.md.

\end{document}
